% Provenance (never printed): see sources/outline_sources.md (Year 2: eccentricity, obliquity and
% precession; insolation computed from spherical astronomy; the Quaternary cycles; carbon dioxide
% and feedbacks). Orbital elements and tabulated insolation of Berger and Loutre (1991) via NOAA
% NCEI (data/derived/bachelor-2-orbit91, fetch_data Book 4 block); periods after Laskar et al.
% (2004) and the NASA Milankovitch page; history after the NASA Earth Observatory feature (2000);
% terminations after Barker and Knorr (2021); the mid-Pleistocene transition after Clark et al.
% (2006) and Koehler and van de Wal (2020). Figures from figdata/bachelor-2/orbital.py; values of
% the exercises and problem computed in the Book 4 scratch calculator calc/ch25.py.
\chapter{कक्षीय प्रणोदन और हिमयुग}\label{ch:b2:orbital-forcing}

पिछले दस लाख वर्षों में पृथ्वी की कक्षा की आकृति ने पूरे ग्रह द्वारा एक वर्ष में प्राप्त सूर्य-प्रकाश को
प्रति हज़ार दो भागों से भी कम बदला है: वायुमंडल के शीर्ष पर \num{340.00} और
\qty{340.54}{W.m^{-2}} के बीच। उन्हीं दस लाख वर्षों में \ang{65} उ पर मध्य-जुलाई का दैनिक-माध्य
सूर्य-प्रकाश लगभग 388 और \qty{500}{W.m^{-2}} के बीच झूलता रहा है, एक-चौथाई से भी
अधिक। हिमयुगों ने दूसरी संख्या का अनुसरण किया, पहली का नहीं। यह अध्याय दोनों को
पृथ्वी की कक्षा के अवयवों से परिकलित करता है, और पूछता है कि एक अक्षांश पर एक ऋतु का सूर्य-प्रकाश
महाद्वीपों को ढकने वाली हिमचादरों की वृद्धि और क्षय की गति कैसे तय कर सकता है।
% ledger: data:b4-orbital, berger91:readme

\begin{recall}
\textbf{भौतिकी की पुस्तकों से}: दीर्घवृत्त, उसका अर्ध-दीर्घ अक्ष $a$ और उत्केंद्रता $e$, केप्लर के
नियम। \textbf{पुस्तक 3 (स्नातक वर्ष 1) से}: क्रांतिवृत्त और तिर्यकता $\varepsilon$, सूर्य की
क्रांति $\delta$, जहाँ $\sin\delta = \sin\varepsilon\sin\lambda$ ($\lambda$ सूर्य का वास्तविक
देशांतर, क्रांतिवृत्त के अनुदिश वसंत विषुव से गिना गया), उष्णकटिबंधीय वर्ष, दैनिक-माध्य
सूर्यातप
\[
  \bar Q = \frac{S_0}{\pi}\left(\frac{a}{r}\right)^{2}\bigl(h_0\sin\varphi\sin\delta +
  \cos\varphi\cos\delta\sin h_0\bigr),\qquad \cos h_0 = -\tan\varphi\tan\delta ,
\]
और उसका वैश्विक और वार्षिक माध्य, $S_0/(4\sqrt{1 - e^2})$। \textbf{स्कूल की पुस्तकों से}:
हिमनदीय चक्र और अंतिम हिमनदीय अधिकतम; विकिरणी प्रणोदन। \textbf{पिछले अध्यायों से}:
\omterm{def:b2:precambrian-earth:ice-albedo}{हिम--एल्बिडो प्रतिपुष्टि} (\cref{def:b2:precambrian-earth:ice-albedo}); नितलवासी स्टैक, EPICA
डोम~C अभिलेख और \omterm{def:b2:climate-proxies:age-model}{आयु-मॉडल} (\cref{ch:b2:climate-proxies})।
\end{recall}
% ledger: earth:TSI, earth:obliquity, earth:ecc

\section{कक्षा और उसके तीन अवयव}\label{sec:b2:orbital-forcing:elements}

पृथ्वी की कक्षा एक दीर्घवृत्त है जिसकी एक नाभि पर सूर्य है। ग्रह, सबसे बढ़कर बृहस्पति और शनि,
उसे धीरे-धीरे विकृत करते हैं और उसके तल में घुमाते हैं; विषुवतीय उभार पर सूर्य और चंद्रमा का खिंचाव
घूर्णन-अक्ष को लट्टू की तरह डगमगाता है। सूर्य-प्रकाश के लिए जो मायने रखता है उसे तीन राशियाँ
बताती हैं: दीर्घवृत्त कितना लंबा है, अक्ष कितना झुका है, और कक्षा के अनुदिश
ऋतुएँ कहाँ पड़ती हैं।
% ledger: nasa:milank-causes

\begin{definition}[उपसौर का देशांतर; जलवायवीय पुरस्सरण]\label{def:b2:orbital-forcing:elements}
\emph{उपसौर का देशांतर}\index{उपसौर का देशांतर} $\varpi$ वह कोण है, जो कक्षा के तल में
वसंत विषुव की दिशा से मापा जाता है, उस दिशा और सूर्य से देखी गई उपसौर की
दिशा के बीच; वह गतिमान विषुव से गिना जाता है, इसलिए वह अक्ष की डगमगाहट और कक्षा के
घूमने, दोनों के साथ बदलता है। \emph{जलवायवीय
पुरस्सरण}\index{जलवायवीय पुरस्सरण} राशि $e\sin\varpi$ है, जिसे पुरस्सरण
सूचकांक भी कहा जाता है।
\end{definition}
% ledger: berger91:readme

\begin{omfigure}
\centering
\begin{tikzpicture}[scale=0.82, font=\scriptsize, >=Stealth]
\draw[thick, rotate=101.37] (-0.72,0) ellipse (2.4 and 2.2895);
\draw[densely dashed, black!60] (-0.331,1.647) -- (0.615,-3.059);
\draw[->, black!70] (0,0) -- (3.9,0);
\node[anchor=south, align=center] at (3.75,0.06) {$\gamma$, वसंत विषुव\\की ओर};
\draw[->, omDef] (0.75,0) arc[start angle=0, end angle=101.37, radius=0.75];
\node[omDef] at (0.78,0.62) {$\varpi$};
\fill[ocRed] (0,0) circle (0.13);
\node[anchor=north west] at (0.05,-0.08) {सूर्य};
\foreach \x/\y in {-2.062/0, 0/-3.094, 2.321/0, 0/1.688} {\fill[ocBlue] (\x,\y) circle (0.08);}
\node[anchor=east, align=right] at (-2.15,0) {मार्च\\विषुव};
\node[anchor=north west, align=left] at (2.38,-0.08) {सितंबर\\विषुव};
\node[anchor=west] at (0.1,1.92) {दिसंबर अयनांत};
\node[anchor=east] at (-0.12,-3.25) {जून अयनांत};
\fill[black] (-0.331,1.647) circle (0.06);
\node[anchor=east] at (-0.42,1.72) {उपसौर};
\fill[black] (0.615,-3.059) circle (0.06);
\node[anchor=west] at (0.7,-3.15) {अपसौर};
\draw[->, thick, black!60] (-1.9,-1.75) arc[start angle=220, end angle=245, radius=2.6];
\end{tikzpicture}
\par\smallskip{\small उत्तर की ओर से देखी गई पृथ्वी की कक्षा, उत्केंद्रता को बढ़ाकर
0.3 किया गया है और उपसौर का आज का देशांतर, $\varpi = \ang{101.4}$। पृथ्वी वामावर्त चलती है
(तीर); बिंदु विषुवों और अयनांतों पर उसकी स्थिति अंकित करते हैं। उपसौर दिसंबर अयनांत के कुछ
दिन बाद पड़ता है, अपसौर जून अयनांत के कुछ दिन बाद। योजनामूलक।}
\end{omfigure}
% ledger: berger91:rows, nasa:milank-detail

पृथ्वी से देखने पर सूर्य क्रांतिवृत्त के अनुदिश चलता है, और उसका वास्तविक देशांतर $\lambda$ वसंत विषुव से
कोण है: जून अयनांत पर $\lambda = \ang{90}$, दिसंबर अयनांत पर \ang{270}। पृथ्वी उपसौर पर तब होती है
जब सूर्य उपसौर की विपरीत दिशा में हो,
$\lambda = \varpi + \ang{180}$ पर; आज $\varpi = \ang{101.37}$ है और उपसौर
$\lambda = \ang{281.37}$ पर पड़ता है, दिसंबर अयनांत के लगभग ग्यारह दिन बाद, 3 जनवरी को या उसके आसपास।
% ledger: berger91:rows, nasa:milank-detail

\begin{definition}[मिलांकोविच प्राचल]\label{def:b2:orbital-forcing:milankovitch}
\emph{मिलांकोविच प्राचल}\index{मिलांकोविच प्राचल} वे तीन कक्षीय
राशियाँ हैं जो ऋतुओं और अक्षांशों में पृथ्वी पर सूर्य-प्रकाश का वितरण तय करती हैं:
उत्केंद्रता $e$, तिर्यकता $\varepsilon$ और \omterm{def:b2:orbital-forcing:elements}{जलवायवीय पुरस्सरण}
$e\sin\varpi$। \emph{मिलांकोविच सिद्धांत}\index{मिलांकोविच सिद्धांत} कहता है कि उनके धीमे
परिवर्तन उच्च उत्तरी अक्षांशों पर ग्रीष्म के सूर्यातप के माध्यम से क्वाटरनरी हिमचादरों की वृद्धि और क्षय की
गति तय करते हैं।
\end{definition}
% ledger: nasa:milank-detail, eo00:milankovitch

हर प्राचल अपनी लयों के साथ बदलता है। उत्केंद्रता के मुख्य आवर्तकाल \num{95000},
\num{99000}, \num{124000} और \num{131000} वर्षों के निकट हैं, जो मिलकर लगभग
\num{100000} वर्षों का एक चक्र बनाते हैं, और \num{405000} वर्षों का एक आवर्तकाल, जो सबमें सबसे स्थिर है; पिछले
दस लाख वर्षों में वह \num{0.0046} से \num{0.0565} तक रही, और आज, \num{0.0172} पर, वह धीरे-धीरे
घट रही है। तिर्यकता \num{40978} वर्षों के मुख्य आवर्तकाल के साथ, पिछले दस लाख वर्षों में \ang{22.1} और
\ang{24.5} के बीच बदलती है; \ang{23.44} पर वह घट रही है, अपने पिछले अधिकतम को पार करके, जो
लगभग \num{10000} वर्ष पहले था। अक्ष लगभग \num{25770} वर्षों में एक चक्कर डगमगाता है, जबकि
कक्षा लगभग \num{112000} में घूमती है; मिलकर वे उपसौर को ऋतुओं में से होकर खिसकाते हैं, जिसके
मुख्य आवर्तकाल \num{19000}, \num{22400} और \num{23700} वर्षों के निकट हैं।
% ledger: la04:eccentricity, la04:obliquity, la04:precession, nasa:milank-detail, data:b4-orbital

\begin{omfigure}
\centering
% figdata/bachelor-2/orbital.py part elements
\begin{tikzpicture}
\begin{groupplot}[group style={group size=1 by 3, vertical sep=0.3cm, xlabels at=edge bottom,
  xticklabels at=edge bottom}, width=0.96\linewidth, height=3.1cm, xmin=0, xmax=1000,
  xtick={0,100,...,1000}, tick label style={font=\scriptsize}, label style={font=\small},
  scaled y ticks=false, yticklabel style={/pgf/number format/fixed,
  /pgf/number format/precision=2}]
\nextgroupplot[ylabel={$e$}, ymin=0, ymax=0.085, ytick={0,0.02,0.04,0.06}]
\addplot[omDef, thin] table[x=age_ka, y=ecc] {figdata/out/bachelor-2-orbital-elements.dat};
\node[font=\scriptsize, anchor=north west] at (rel axis cs:0.01,0.97) {क\enspace उत्केंद्रता};
\nextgroupplot[ylabel={$\varepsilon$ (\unit{\degree})}, ymin=21.9, ymax=25.4, ytick={22,23,24}]
\addplot[ocBlue, thin] table[x=age_ka, y=obl] {figdata/out/bachelor-2-orbital-elements.dat};
\node[font=\scriptsize, anchor=north west] at (rel axis cs:0.01,0.97) {ख\enspace तिर्यकता};
\nextgroupplot[ylabel={$e\sin\varpi$}, ymin=-0.07, ymax=0.12, ytick={-0.05,0,0.05},
  xlabel={आयु (1950 से पहले हज़ार वर्ष)}]
\addplot[ocRed, thin] table[x=age_ka, y=prec] {figdata/out/bachelor-2-orbital-elements.dat};
\node[font=\scriptsize, anchor=north west] at (rel axis cs:0.01,0.97) {ग\enspace जलवायवीय
पुरस्सरण};
\end{groupplot}
\end{tikzpicture}
\par\smallskip{\small पिछले दस लाख वर्षों में तीन \omterm{def:b2:orbital-forcing:milankovitch}{मिलांकोविच प्राचल}: (क)
उत्केंद्रता; (ख) तिर्यकता; (ग) \omterm{def:b2:orbital-forcing:elements}{जलवायवीय पुरस्सरण}, जिसका आवरण उत्केंद्रता है।
डेटा: NOAA NCEI Paleoclimatology, Berger और Loutre (1991), 23 अक्टूबर 2026 को प्राप्त।}
\end{omfigure}
% ledger: data:b4-orbit91

\omterm{def:b2:orbital-forcing:elements}{जलवायवीय पुरस्सरण} तेज़ी से झूलता है, पर केवल उतना जितना उत्केंद्रता उसे करने देती है: जब कक्षा
लगभग वृत्ताकार हो, जैसे \num{400000} वर्ष पहले, तो इससे शायद ही कोई अंतर पड़ता है कि उपसौर कहाँ पड़ता है।
चित्र इसे स्पष्ट दिखाता है: पैनल (ग) का आवरण पैनल (क) है।
% ledger: data:b4-orbit91

\section{कक्षा से सूर्य-प्रकाश}\label{sec:b2:orbital-forcing:insolation}

ऊपर स्मरण किए गए दैनिक-माध्य सूर्यातप को एक और घटक चाहिए: किसी दिए गए दिन सूर्य की
दूरी।

\begin{proposition}[दूरी और सूर्य का देशांतर]\label{prop:b2:orbital-forcing:distance}
जब सूर्य का वास्तविक देशांतर $\lambda$ हो, तो पृथ्वी--सूर्य दूरी इस प्रकार है:
\[
  \frac{r}{a} = \frac{1 - e^2}{1 - e\cos(\lambda - \varpi)} .
\]
जून अयनांत पर ($\lambda = \ang{90}$), $e$ में प्रथम कोटि तक,
$(a/r)^2 \approx 1 - 2e\sin\varpi$; दिसंबर अयनांत पर, $(a/r)^2 \approx 1 + 2e\sin\varpi$।
\end{proposition}

\begin{proof}
सूर्य को नाभि पर रखने वाले दीर्घवृत्त पर, $r/a = (1 - e^2)/(1 + e\cos v)$, जहाँ वास्तविक विसंगति $v$
सूर्य पर उपसौर से पृथ्वी तक का कोण है। सूर्य से पृथ्वी उस दिशा के विपरीत दिशा में
दिखती है जिसमें पृथ्वी से सूर्य दिखता है, इसलिए सूर्य से देखा गया उसका देशांतर
$\lambda + \ang{180}$ है, और $v = \lambda + \ang{180} - \varpi$, $\cos v = -\cos(\lambda - \varpi)$।
$\lambda = \ang{90}$ पर, $\cos(\lambda - \varpi) = \sin\varpi$ और $(a/r)^2 = (1 -
e\sin\varpi)^2/(1 - e^2)^2 = 1 - 2e\sin\varpi + O(e^2)$; $\lambda = \ang{270}$ पर चिह्न बदल जाता है।
\end{proof}

इसलिए \omterm{def:b2:orbital-forcing:elements}{जलवायवीय पुरस्सरण} ही सही सूचकांक है: वह मापता है कि जून अयनांत पर पृथ्वी सूर्य के कितनी अधिक निकट,
या कितनी अधिक दूर, है। जब $e\sin\varpi$ बड़ा और धनात्मक हो, जैसे आज और
\num{116000} वर्ष पहले कहीं अधिक, तो उत्तरी ग्रीष्म अपसौर के निकट पड़ता है और दुर्बल होता है; जब वह
बड़ा और ऋणात्मक हो, जैसे लगभग \num{10000} वर्ष पहले, तो उत्तरी ग्रीष्म उपसौर के निकट पड़ता है।
% ledger: berger91:rows

\begin{example}[\ang{65} उ पर जून का सूर्य, आज]\label{ex:b2:orbital-forcing:today}
आज $e = 0.017236$, $\varpi = \ang{101.37}$, $\varepsilon = \ang{23.446}$, और तालिकाएँ
$S_0 = \qty{1360}{W.m^{-2}}$ का उपयोग करती हैं। जून अयनांत पर $\delta = \varepsilon$; $\varphi =
\ang{65}$ पर, $\cos h_0 = -\tan\ang{65}\tan\ang{23.446} = -0.9301$, अतः $h_0 = \ang{158.44} =
\qty{2.7654}{rad}$: सूर्य $2 \times 158.44/360 \times 24 = 21.1$~घंटे ऊपर रहता है।
सूर्यातप सूत्र का कोष्ठक $2.7654 \times 0.9063 \times 0.3979 + 0.4226 \times 0.9174 \times 0.3674 =
1.1397$ है, और माध्य दूरी पर $\bar Q = 432.90 \times 1.1397 = \qty{493.4}{W.m^{-2}}$। दूरी
गुणक $(a/r)^2 = 0.9671$ है, अतः $\bar Q = \qty{477.1}{W.m^{-2}}$।
\end{example}
% ledger: berger91:rows, berger91:readme, data:b4-orbital

केप्लर का दूसरा नियम एक मोड़ जोड़ता है। अपसौर के निकट पृथ्वी अधिक धीरे चलती है: वहाँ पड़ने वाली ऋतु
दुर्बल होती है पर अधिक लंबी चलती है। उत्तरी ग्रीष्म आज उत्तरी शीत ऋतु से लगभग 4.5 दिन लंबा है,
और वसंत शरद से लगभग तीन दिन लंबा।
% ledger: nasa:milank-detail

\begin{proposition}[वार्षिक माध्य पुरस्सरण पर निर्भर नहीं करते]\label{prop:b2:orbital-forcing:annual}
सूर्य के दो दिए गए देशांतरों के बीच वायुमंडल के शीर्ष पर प्रति इकाई क्षेत्रफल प्राप्त ऊर्जा
$\varpi$ पर निर्भर नहीं करती, और किसी भी अक्षांश पर वार्षिक-माध्य सूर्यातप त्रिज्या $a$ की
वृत्ताकार कक्षा के सूर्यातप का $1/\sqrt{1 - e^2}$ गुना है। किसी ध्रुव पर वार्षिक माध्य यह है:
\[
  \bar Q_{\text{ध्रुव}} = \frac{S_0\sin\varepsilon}{\pi\sqrt{1 - e^2}} ,
\]
जो तिर्यकता के साथ बढ़ता है।
\end{proposition}

\begin{proof}
केप्लर के दूसरे नियम से त्रिज्या स्थिर दर $\pi ab/T$ से क्षेत्रफल बुहारती है, और
पृथ्वी के $\mathrm d\lambda$ घूमने के दौरान बुहारा गया क्षेत्रफल $\frac12 r^2\,\mathrm d\lambda$ है, अतः $\mathrm dt =
T r^2\,\mathrm d\lambda/(2\pi ab)$, जहाँ $b = a\sqrt{1 - e^2}$। किसी दिए गए अक्षांश
और सूर्य के देशांतर पर सूर्यातप $(a/r)^2$ गुणा उसका मान $Q_a(\lambda)$ है, जो दूरी $a$ पर है, इसलिए
प्राप्त ऊर्जा $\int Q_a(\lambda)(a/r)^2\,\mathrm dt = \frac{T}{2\pi\sqrt{1 - e^2}}\int
Q_a(\lambda)\,\mathrm d\lambda$ है: दूरी निरस्त हो गई है, और उसके साथ $\varpi$। पूरे
वर्ष में $\lambda$ \ang{360} से होकर गुज़रता है और माध्य $\frac{1}{2\pi\sqrt{1-e^2}}\int_0^{2\pi}
Q_a\,\mathrm d\lambda$ है। उत्तरी ध्रुव पर $Q_a = S_0\sin\delta = S_0\sin\varepsilon\sin\lambda$ है, $0 < \lambda < \pi$ के लिए,
और अन्यथा 0; समाकल $2S_0\sin\varepsilon$ है।
\end{proof}
% ledger: earth:TSI

पुरस्सरण सूर्य-प्रकाश को एक ऋतु से दूसरी ऋतु में ले जाता है, कभी किसी वर्ष के भीतर या बाहर नहीं; उत्केंद्रता
वार्षिक योग को $1/\sqrt{1 - e^2}$ गुणक से बदलती है, पिछले दस लाख वर्षों में अधिकतम \qty{0.16}{\percent};
तिर्यकता सूर्य-प्रकाश को अक्षांशों के बीच ले जाती है, ध्रुवों पर वार्षिक माध्य बढ़ाती है और
विषुवत रेखा पर उसे थोड़ा घटाती है। केवल ऋतुवार और अक्षांशीय प्रतिरूप
बहुत बदलता है, और नीचे का चित्र यही दिखाता है।
% ledger: data:b4-orbital, nasa:milank-detail

\begin{omfigure}
\centering
% figdata/bachelor-2/orbital.py part lat
\begin{tikzpicture}
\begin{groupplot}[group style={group size=2 by 1, horizontal sep=1.6cm},
  tick label style={font=\scriptsize}, label style={font=\small}, height=5.4cm,
  width=0.48\linewidth, xmin=-90, xmax=90, xtick={-90,-60,...,90}, xlabel={अक्षांश
  (\unit{\degree})}, ylabel={$\bar Q$ (\unit{W.m^{-2}})}]
\nextgroupplot[ymin=0, ymax=650, title={क\enspace जून अयनांत}, title style={font=\scriptsize},
  legend style={font=\scriptsize, at={(0.03,0.97)}, anchor=north west, draw=none},
  legend cell align=left]
\addplot[black, thick] table[x=lat, y=jun_0] {figdata/out/bachelor-2-orbital-lat.dat};
\addlegendentry{आज}
\addplot[ocRed, thick, densely dashed] table[x=lat, y=jun_10] {figdata/out/bachelor-2-orbital-lat.dat};
\addlegendentry{\num{10000} वर्ष पहले}
\nextgroupplot[ymin=100, ymax=500, title={ख\enspace वार्षिक माध्य}, title style={font=\scriptsize},
  legend style={font=\scriptsize, at={(0.5,0.03)}, anchor=south, draw=none},
  legend cell align=left]
\addplot[ocBlue, thick] table[x=lat, y=ann_lo] {figdata/out/bachelor-2-orbital-lat.dat};
\addlegendentry{$\varepsilon = \ang{22.1}$}
\addplot[ocRed, thick, densely dashed] table[x=lat, y=ann_hi] {figdata/out/bachelor-2-orbital-lat.dat};
\addlegendentry{$\varepsilon = \ang{24.5}$}
\end{groupplot}
\end{tikzpicture}
\par\smallskip{\small (क) जून अयनांत पर दैनिक-माध्य सूर्यातप, आज और \num{10000} वर्ष
पहले, जब उत्तरी ग्रीष्म उपसौर के निकट पड़ता था और तिर्यकता अधिक थी: उच्च उत्तरी अक्षांशों में लगभग
\qty{50}{W.m^{-2}} अधिक। (ख) आज की उत्केंद्रता के साथ, पिछले दस लाख वर्षों की दो
चरम तिर्यकताओं के लिए वार्षिक-माध्य सूर्यातप: अधिक झुकाव ध्रुवों को
अधिक और उष्णकटिबंध को थोड़ा कम सूर्य-प्रकाश भेजता है। कक्षीय अवयवों से
$S_0 = \qty{1360}{W.m^{-2}}$ के साथ परिकलित।}
\end{omfigure}
% ledger: data:b4-orbital, data:b4-orbit91

\section{सूर्यातप वक्र का परिकलन}\label{sec:b2:orbital-forcing:curve}

मिलांकोविच ने जो सूर्यातप वक्र बनाया था, वह अब एक नियमित परिकलन है। अतीत के कक्षीय अवयव
ग्रहों की गतियों का समाकलन करके परिकलित किए जाते हैं; हर हज़ार वर्ष पर सारणीबद्ध,
वे निवेश हैं। हर आयु के लिए नीचे की विधि किसी भी अक्षांश और ऋतु पर सूर्यातप देती है।

\begin{method}[अतीत का सूर्यातप]\label{met:b2:orbital-forcing:curve}
\begin{enumerate}
  \item वांछित आयु के लिए $e$, $\varpi$ और $\varepsilon$ पढ़िए।
  \item सूर्य के वास्तविक देशांतर $\lambda$ से ऋतु चुनिए: जून अयनांत के लिए \ang{90};
        तालिकाओं का मध्य-जुलाई $\lambda = \ang{120}$ के संगत है।
  \item $\delta$ परिकलित कीजिए, $\sin\delta = \sin\varepsilon\sin\lambda$ से, और सूर्यास्त का घंटा कोण $h_0$,
        अक्षांश $\varphi$ पर (यदि $|\tan\varphi\tan\delta| > 1$ हो तो ध्रुवीय दिन या रात)।
  \item \cref{prop:b2:orbital-forcing:distance} से दूरी गुणक $(a/r)^2$ परिकलित कीजिए।
  \item गुणा कीजिए: $\bar Q = (S_0/\pi)(a/r)^2(h_0\sin\varphi\sin\delta +
        \cos\varphi\cos\delta\sin h_0)$।
\end{enumerate}
\end{method}
% ledger: berger91:readme, data:b4-orbital

पिछले दस लाख वर्षों की 1001 आयुओं पर लागू करने पर, यह विधि \ang{65} उ पर सारणीबद्ध मध्य-जुलाई
सूर्यातप को \qty{0.012}{W.m^{-2}} के भीतर, यानी तालिकाओं के पूर्णांकन तक, पुनरुत्पन्न करती है, और
वर्तमान दिन के सारणीबद्ध जून मानों को ठीक-ठीक। तालिकाएँ समय का मूल 1950 पर रखती हैं।
% ledger: data:b4-orbital, berger91:readme

\begin{omfigure}
\centering
% figdata/bachelor-2/orbital.py part jul65
\begin{tikzpicture}
\begin{groupplot}[group style={group size=1 by 2, vertical sep=0.3cm, xlabels at=edge bottom,
  xticklabels at=edge bottom}, width=0.96\linewidth, xmin=0, xmax=400, xtick={0,50,...,400},
  tick label style={font=\scriptsize}, label style={font=\small}]
\nextgroupplot[height=5.0cm, ymin=370, ymax=590, ylabel={$\bar Q$, \ang{65} उ पर
  (\unit{W.m^{-2}})}, legend style={font=\scriptsize, at={(0.5,1.02)}, anchor=south, draw=none,
  legend columns=3, /tikz/every even column/.append style={column sep=6pt}},
  legend cell align=left]
\addplot[black!45, line width=2.2pt] table[x=age_ka, y=tab] {figdata/out/bachelor-2-orbital-jul65.dat};
\addlegendentry{मध्य-जुलाई, सारणीबद्ध}
\addplot[ocRed, thin] table[x=age_ka, y=calc] {figdata/out/bachelor-2-orbital-jul65.dat};
\addlegendentry{मध्य-जुलाई, परिकलित}
\addplot[ocBlue, thin, densely dashed] table[x=age_ka, y=jun65] {figdata/out/bachelor-2-orbital-jul65.dat};
\addlegendentry{जून अयनांत, परिकलित}
\nextgroupplot[height=2.8cm, ymin=-0.02, ymax=0.02, ytick={-0.01,0,0.01},
  scaled y ticks=false, yticklabel style={/pgf/number format/fixed}, ylabel={अवशेष},
  xlabel={आयु (1950 से पहले हज़ार वर्ष)}]
\addplot[only marks, mark=*, mark size=0.35pt, black] table[x=age_ka, y=resid]
  {figdata/out/bachelor-2-orbital-jul65.dat};
\end{groupplot}
\end{tikzpicture}
\par\smallskip{\small ऊपर: \ang{65} उ पर दैनिक-माध्य सूर्यातप, पिछले \num{400000} वर्षों में,
मध्य-जुलाई का, जैसा बर्जर और लूत्र ने सारणीबद्ध किया (धूसर, चौड़ा) और जैसा ऊपर की विधि से परिकलित (लाल),
और जून अयनांत पर (नीला, धराशायी)। नीचे: परिकलित घटा सारणीबद्ध मध्य-जुलाई मान,
\unit{W.m^{-2}} में।}
\end{omfigure}
% ledger: data:b4-orbit91, data:b4-orbital

दो विशेषताएँ उभरती हैं। वक्र पर \omterm{def:b2:orbital-forcing:elements}{जलवायवीय पुरस्सरण} हावी है, जिसके शिखर लगभग
\num{20000} वर्ष के अंतर पर हैं और जिनकी ऊँचाई उत्केंद्रता के साथ फूलती और सिकुड़ती है; तिर्यकता
एक धीमा मॉडुलन जोड़ती है। और उतार-चढ़ाव बड़े हैं: पिछले \num{800000} वर्षों में मध्य-जुलाई का मान
388 से \qty{497}{W.m^{-2}} तक रहा।
% ledger: data:b4-orbital

\section{मिलांकोविच सिद्धांत और हिमनदीय चक्र}\label{sec:b2:orbital-forcing:theory}

\begin{history}[पैंसठ अंश उत्तर पर धूप]
1879 में जन्मे मिलुटिन मिलांकोविच ने सूर्य-प्रकाश के ऋतुवार और अक्षांशीय परिवर्तनों से जलवायु का एक गणितीय सिद्धांत बनाया। तीनों कक्षीय परिवर्तनों को मिलाकर,
उन्होंने 1800 से पहले के \num{600000} वर्षों के लिए अक्षांशों का सूर्यातप और संगत सतही तापमान
परिकलित किए। जलवायुविद कोपेन के सुझाव पर उन्होंने \ang{65} उ पर ग्रीष्म के
सूर्यातप को निर्णायक राशि चुना: विशाल हिमचादरें उसी अक्षांश के निकट
बढ़ी थीं, और ठंडी ग्रीष्म ऋतुएँ शीत ऋतु की हिम को कम पिघलाती हैं, जिससे हिम वर्ष-दर-वर्ष
संचित होती है। वे तिर्यकता को तीनों प्राचलों में सबसे महत्त्वपूर्ण मानते थे। लगभग पचास वर्षों तक
यह सिद्धांत बड़े पैमाने पर उपेक्षित रहा; मिलांकोविच की मृत्यु 1958 में हुई। 1976 में हेज़, इम्ब्री और शैकलटन
ने पिछले \num{450000} वर्षों के गहरे समुद्र के अवसादों में जलवायु परिवर्तिता के स्पेक्ट्रमी शिखर
लगभग \num{23000}, \num{42000} और \num{100000} वर्षों पर पाए, जो पुरस्सरण, तिर्यकता और
उत्केंद्रता के आवर्तकाल हैं, और निष्कर्ष निकाला कि पृथ्वी की कक्षा में परिवर्तन क्वाटरनरी हिमयुगों का
मूलभूत कारण हैं।
\end{history}
% ledger: eo00:milankovitch, nasa:milank-detail, hays76:periods

नीचे का चित्र \cref{ch:b2:climate-proxies} के अभिलेखों पर उनका परीक्षण दोहराता है। वह \ang{65} उ के
जुलाई सूर्यातप और नितलवासी स्टैक के आयाम स्पेक्ट्रमों, यानी हर आवर्तकाल की प्रबलता, की तुलना करता है,
\num{700000} वर्षों की दो खिड़कियों में।

\begin{omfigure}
\centering
% figdata/bachelor-2/orbital.py part spectrum
\begin{tikzpicture}
\begin{axis}[width=0.9\linewidth, height=5.6cm, xmin=10, xmax=150, ymin=0, ymax=1.08,
  xlabel={आवर्तकाल (हज़ार वर्ष)}, ylabel={सापेक्ष आयाम},
  xtick={19,23,41,100,150}, tick label style={font=\scriptsize}, label style={font=\small},
  legend style={font=\scriptsize, at={(0.5,1.02)}, anchor=south, draw=none, legend columns=3,
  /tikz/every even column/.append style={column sep=6pt}}, legend cell align=left]
\addplot[black!50, thick] table[x=period, y=amp_ins] {figdata/out/bachelor-2-orbital-spectrum.dat};
\addlegendentry{सूर्यातप \ang{65} उ, 0--\qty{700}{ka}}
\addplot[ocRed, thick] table[x=period, y=amp_young] {figdata/out/bachelor-2-orbital-spectrum.dat};
\addlegendentry{LR04, 0--\qty{700}{ka}}
\addplot[ocBlue, thick, densely dashed] table[x=period, y=amp_old] {figdata/out/bachelor-2-orbital-spectrum.dat};
\addlegendentry{LR04, \numrange{1250}{1950}~ka}
\end{axis}
\end{tikzpicture}
\par\smallskip{\small आयाम स्पेक्ट्रम, हर एक अपने अधिकतम के अनुसार मापित: \ang{65} उ पर
मध्य-जुलाई सूर्यातप का, पिछले \num{700000} वर्षों में, और LR04 नितलवासी स्टैक का, उसी खिड़की में
और \num{1250000} से \num{1950000} वर्ष पहले तक (प्रवृत्ति हटाई गई, टेपर लगाया गया)। सूर्यातप की
शक्ति पुरस्सरण के आवर्तकालों पर, 19 से 24 हज़ार वर्षों के निकट, और तिर्यकता के आवर्तकाल पर है; हाल का
हिम-आयतन \num{100000} वर्षों के निकट शिखर पर है, जहाँ सूर्यातप में लगभग कुछ नहीं है; पुराना
हिम-आयतन \num{41000} वर्षों पर शिखर पर है।}
\end{omfigure}
% ledger: data:b4-orbital, data:lr04, hays76:periods

\subsection*{हिमनदीय काल, अंतर्हिमनदीय काल और समापन}

\begin{definition}[हिमनदीय काल; अंतर्हिमनदीय काल; हिमनदीय समापन]\label{def:b2:orbital-forcing:glacials}
क्वाटरनरी अभिलेख में \emph{हिमनदीय काल}\index{हिमनदीय काल} वह अंतराल है जब विशाल
हिमचादरें उत्तरी महाद्वीपों को ढकती हैं, जो उच्च नितलवासी $\delta^{18}$O और निम्न समुद्र-तल से दर्ज होता है;
\emph{अंतर्हिमनदीय काल}\index{अंतर्हिमनदीय काल} वह अंतराल है जिसमें ग्रीनलैंड के बाहर उत्तरी
गोलार्ध में स्थल पर बहुत कम बर्फ़ होती है, वर्तमान की तरह। \emph{हिमनदीय
समापन}\index{हिमनदीय समापन} उन प्रमुख संक्रमणों में से एक है जो किसी हिमनदीय चक्र को समाप्त करते हैं,
पूर्ण हिमनदीय से अंतर्हिमनदीय परिस्थितियों तक, जो अपने से पहले की वृद्धि से कहीं अधिक तेज़ होता है।
\end{definition}
% ledger: kohler20:mpt, barker21:terminations, ipcc:ref-periods

पिछले पाँच लाख वर्षों में जलवायु का प्रभावी आवर्तकाल \num{100000} वर्षों के निकट है,
यद्यपि उस आवर्तकाल पर प्रणोदन, जो केवल उत्केंद्रता से आता है, तिर्यकता और पुरस्सरण के प्रणोदन
की तुलना में नगण्य है। चक्र एकतरफ़ा भी हैं: नितलवासी $\delta^{18}$O
लगभग \num{90000} वर्षों में बर्फ़ का सीढ़ीनुमा संचय दर्ज करता है, फिर लगभग
\num{10000} में क्षय, एक आरी-दाँत आकृति जिसके गिरते किनारे समापन हैं। अंतिम \omterm{def:b2:orbital-forcing:glacials}{अंतर्हिमनदीय काल}, लगभग
\num{129000} से \num{116000} वर्ष पहले तक, तब समाप्त हुआ जब \ang{65} उ पर ग्रीष्म का सूर्यातप गिरकर
पिछले \num{800000} वर्षों के अपने सबसे निम्न मानों में से एक पर पहुँचा: जून-अयनांत मान अपने
न्यूनतम, \qty{438.6}{W.m^{-2}}, पर \num{116000} वर्ष पहले पहुँचा, और मध्य-जुलाई मान अपने दूसरे सबसे निम्न
न्यूनतम, \qty{392.2}{W.m^{-2}}, पर \num{114000} वर्ष पहले। अंतिम हिमनदीय अधिकतम लगभग
\num{95000} वर्ष बाद आया।
% ledger: barker21:terminations, ipcc:ref-periods, data:b4-orbital, data:b4-orbital-extremes

कोई हिमचादर अपने प्रणोदन का कितनी जल्दी अनुसरण कर सकती है? सबसे सरल मॉडल हिम-आयतन $V$ को
उस आयतन $V_{\mathrm{eq}}$ की ओर शिथिल होता हुआ मानता है जिसे वर्तमान सूर्यातप बनाए रखेगा, एक
अनुक्रिया काल $\tau$ के साथ:
\[
  \frac{\mathrm dV}{\mathrm dt} = \frac{V_{\mathrm{eq}}(t) - V}{\tau} .
\]

\begin{proposition}[एक पश्चित अनुक्रिया]\label{prop:b2:orbital-forcing:lag}
प्रणोदन $V_{\mathrm{eq}} = A\sin\omega t$ के लिए पश्चता मॉडल का हल इसकी ओर प्रवृत्त होता है
\[
  V = \frac{A}{\sqrt{1 + \omega^2\tau^2}}\,\sin\bigl(\omega(t - t_{\ell})\bigr),\qquad
  t_{\ell} = \frac{\arctan(\omega\tau)}{\omega} :
\]
बर्फ़ धीमे प्रणोदन का लगभग पूरा और तेज़ प्रणोदन का केवल दुर्बल रूप से अनुसरण करती है, सदा देर से। ज्यावक्रों के
योग की अनुक्रिया अनुक्रियाओं का योग है: मॉडल कोई ऐसा आवर्तकाल नहीं बनाता जो उसके प्रणोदन में
न हो।
\end{proposition}

\begin{proof}
$V = B\sin\omega t + C\cos\omega t$ आज़माइए: प्रतिस्थापित करने पर $\omega B\tau = -C$ और $-\omega C\tau = A -
B$, अतः $B = A/(1 + \omega^2\tau^2)$, $C = -\omega\tau A/(1 + \omega^2\tau^2)$, आयाम
$\sqrt{B^2 + C^2} = A/\sqrt{1 + \omega^2\tau^2}$ और कला-विलंब $\arctan(\omega\tau)$। कोई भी अन्य
हल इससे $\mathrm e^{-t/\tau}$ के एक गुणज से भिन्न होता है, जो क्षीण होकर मिट जाता है।
समीकरण रैखिक है, इसलिए हल जुड़ते हैं।
\end{proof}

\begin{omfigure}
\centering
% figdata/bachelor-2/orbital.py part lag
\begin{tikzpicture}
\begin{axis}[width=0.96\linewidth, height=5.4cm, xmin=0, xmax=400, ymin=-3.4, ymax=4.2,
  xtick={0,50,...,400}, xlabel={आयु (1950 से पहले हज़ार वर्ष)},
  ylabel={मानकीकृत}, tick label style={font=\scriptsize}, label style={font=\small},
  legend style={font=\scriptsize, at={(0.5,1.02)}, anchor=south, draw=none, legend columns=3,
  /tikz/every even column/.append style={column sep=6pt}}, legend cell align=left]
\addplot[black!35, thin] table[x=age_ka, y=forcing] {figdata/out/bachelor-2-orbital-lag.dat};
\addlegendentry{ऋण सूर्यातप \ang{65} उ}
\addplot[ocRed, thick] table[x=age_ka, y=model] {figdata/out/bachelor-2-orbital-lag.dat};
\addlegendentry{पश्चता मॉडल, $\tau = \qty{10}{kyr}$}
\addplot[black, thick] table[x=age_ka, y=lr04] {figdata/out/bachelor-2-orbital-lag.dat};
\addlegendentry{LR04 (बर्फ़ ऊपर)}
\end{axis}
\end{tikzpicture}
\par\smallskip{\small \ang{65} उ पर मध्य-जुलाई सूर्यातप से चालित पश्चता मॉडल (कम
धूप, अधिक बर्फ़), नितलवासी स्टैक के सापेक्ष, तीनों मानकीकृत (माध्य 0, मानक विचलन
1); धनात्मक मानों का अर्थ अधिक बर्फ़ है। मॉडल स्टैक की पुरस्सरण-लहरियों का अनुसरण करता है पर उसमें
उसके बड़े आरी-दाँत चक्रों में से कोई नहीं है। \qty{10}{kyr} का अनुक्रिया काल एक मॉडल-चयन है।}
\end{omfigure}
% ledger: data:b4-orbital, data:lr04

पिछले \num{800000} वर्षों में पश्चित मॉडल स्टैक से $r = 0.42$ पर सहसंबद्ध है,
स्वयं सूर्यातप केवल $0.23$ पर: पश्चता सहायक है, पर एक रैखिक हिमचादर ऐसे प्रणोदन से
\num{100000}-वर्षीय चक्र नहीं बना सकती जिसमें वे हों ही नहीं। समापनों में
अरैखिक व्यवहार शामिल होना चाहिए जो सूर्यातप के एक मामूली धक्के को प्रवर्धित करे। उम्मीदवार वे प्रतिपुष्टियाँ हैं जो पूरी जलवायु पर कार्य करती हैं:
\cref{def:b2:precambrian-earth:ice-albedo} की \omterm{def:b2:precambrian-earth:ice-albedo}{हिम--एल्बिडो प्रतिपुष्टि}, क्योंकि अधिक
हिम और बर्फ़ वाला ग्रह अधिक सूर्य-प्रकाश परावर्तित करता है; और कार्बन डाइऑक्साइड, जो अंटार्कटिक बर्फ़ के कोरों में
सहस्राब्दियों के समय-पैमानों पर तब-तब बढ़ती है जब अटलांटिक का अधिक्रमण परिसंचरण दुर्बल होता है, जैसा वह
तब होता है जब हिमचादरों का पिघला जल उत्तर अटलांटिक में उँडेला जाता है। कार्बन
डाइऑक्साइड की एक देहली से आगे, अंतर्हिमनदीय परिस्थितियाँ स्थिर हो सकती हैं।
% ledger: data:b4-orbital, barker21:terminations, eo00:milankovitch

\begin{definition}[मध्य-प्लाइस्टोसीन संक्रमण]\label{def:b2:orbital-forcing:mpt}
\emph{मध्य-प्लाइस्टोसीन संक्रमण}\index{मध्य-प्लाइस्टोसीन संक्रमण} लगभग
\num{1250000} से \num{700000} वर्ष पहले तक का वह परिवर्तन है, लगभग \num{41000} वर्षों और छोटे
आयाम के हिमनदीय चक्रों से लगभग \num{100000} वर्षों और बड़े आयाम के अर्ध-आवर्ती चक्रों की ओर, अधिक मोटी
हिमचादरों के साथ, कक्षीय प्रणोदन में किसी संगत परिवर्तन के बिना।
\end{definition}
% ledger: clark06:mpt, nasa:milank-detail, kohler20:mpt

संक्रमण से पहले हिमचादरें तिर्यकता का अनुसरण करती थीं, जैसा ऊपर के चित्र का धराशायी
स्पेक्ट्रम दिखाता है; एक गणना आरंभिक, मध्य और उत्तर क्वाटरनरी के तिर्यकता-चक्रों में से 67, 88 और \qty{52}{\percent}
में एक \omterm{def:b2:orbital-forcing:glacials}{हिमनदीय समापन} पाती है। केवल कक्षीय सिद्धांत इस परिवर्तन को नहीं समझा सकता; उम्मीदवार
स्वयं हिमचादरों की गतिकी हैं, संभवतः उस नरम रेगोलिथ के हटने के बाद जिस पर वे फिसलती थीं,
और हिमनदीय न्यूनतमों के दौरान कार्बन डाइऑक्साइड की क्रमिक कमी। अभी तक कोई व्याख्या
सामान्य रूप से स्वीकृत नहीं है।
% ledger: kohler20:mpt, nasa:milank-detail, clark06:mpt

\section{हिमनदीय चक्रों में कार्बन डाइऑक्साइड}\label{sec:b2:orbital-forcing:co2}

अंटार्कटिक बर्फ़ के कोर दिखाते हैं कि कार्बन डाइऑक्साइड हिमनदीय चक्रों का अनुसरण करती है: अंतिम हिमनदीय अधिकतम पर 188 और
\qty{194}{ppm} के बीच, अंतिम \omterm{def:b2:orbital-forcing:glacials}{अंतर्हिमनदीय काल} में 266 और \qty{282}{ppm} के बीच;
\cref{ch:b2:climate-proxies} का संयुक्त अभिलेख \qty{189.7}{ppm} देता है, 19 से \qty{23}{ka} के लिए,
और उद्योग से पहले के दो हज़ार वर्षों में \qty{279.4}{ppm}। उसका विकिरणी प्रणोदन
सांद्रता के लघुगणक के लगभग अनुसार बढ़ता है, हर दोगुनीकरण पर \qty{3.93}{W.m^{-2}}।
% ledger: ipcc:fig2-34, data:b4-proxies-co2, ipcc:log-co2, ipcc:erf-2xco2

\begin{example}[कक्षा बनाम कार्बन डाइऑक्साइड]\label{ex:b2:orbital-forcing:co2}
\qty{189.7}{ppm} से \qty{279.4}{ppm} तक प्रणोदन, पूरे ग्लोब पर और पूरे वर्ष में, $3.93\log_2(279.4/189.7) = 3.93 \times
0.559 = \qty{2.2}{W.m^{-2}}$ है। कक्षा ने पिछले
दस लाख वर्षों में सूर्यातप के वैश्विक और वार्षिक माध्य को अधिकतम $340.54 - 340.00 =
\qty{0.54}{W.m^{-2}}$ बदला, पर \ang{65} उ पर मध्य-जुलाई सूर्यातप को लगभग \qty{109}{W.m^{-2}}।
\end{example}
% ledger: data:b4-proxies-co2, ipcc:erf-2xco2, ipcc:log-co2, data:b4-orbital

कक्षा एक गति-निर्धारक है, इंजन नहीं। वह सूर्य-प्रकाश को ऋतुओं और अक्षांशों के बीच पुनर्वितरित करती है
और समय तय करती है; हिमचादरें, अपने एल्बिडो और अपनी मंदता से, और कार्बन डाइऑक्साइड, अपने
वैश्विक प्रणोदन से, एक क्षेत्रीय संकेत को जलवायु के एक वैश्विक परिवर्तन में बदल देते हैं।
% ledger: hays76:periods, barker21:terminations, eo00:milankovitch

\section*{अभ्यास}\label{sec:b2:orbital-forcing:exercises}

अभ्यासों के आँकड़े, अन्यथा न कहा जाए तो: $S_0 = \qty{1360}{W.m^{-2}}$ (तालिकाओं का
मान); \num{365.2422} दिनों का उष्णकटिबंधीय वर्ष; तालिकाओं के कक्षीय अवयव, जिनमें आज
$e = 0.017236$, $\varpi = \ang{101.37}$, $\varepsilon = \ang{23.446}$।
% ledger: berger91:readme, berger91:rows

\begin{exercise}[$\star$]\label{exo:b2:orbital-forcing:1}
$e = 0.0167$ लीजिए। (क)~अपसौर दूरी $a(1 + e)$ और उपसौर
दूरी $a(1 - e)$ का अनुपात परिकलित कीजिए। (ख)~उनके अंतर को $a$ के प्रतिशत के रूप में व्यक्त कीजिए। (ग)~उपसौर पर
सूर्य-प्रकाश अपसौर की तुलना में किस गुणक से अधिक प्रबल है?
\end{exercise}
% ledger: earth:ecc, nasa:milank-detail

\begin{exercise}[$\star$]\label{exo:b2:orbital-forcing:2}
तालिकाएँ आज $e = 0.017236$, $\varpi = \ang{101.37}$ देती हैं; $e = 0.019971$, $\varpi =
\ang{293.86}$ \num{10000} वर्ष पहले; $e = 0.019249$, $\varpi = \ang{359.99}$ \num{6000} वर्ष पहले।
(क)~हर आयु के लिए \omterm{def:b2:orbital-forcing:elements}{जलवायवीय पुरस्सरण} परिकलित कीजिए। (ख)~हर एक के लिए उपसौर पर सूर्य का
देशांतर दीजिए। (ग)~किस आयु पर उपसौर जून अयनांत के सबसे निकट था, और किस आयु पर
सितंबर विषुव के?
\end{exercise}
% ledger: berger91:rows

\begin{exercise}[$\star$]\label{exo:b2:orbital-forcing:3}
घूर्णन-अक्ष तारों के सापेक्ष लगभग \num{25770} वर्षों के आवर्तकाल से पुरस्सरण करता है, जबकि
कक्षा लगभग \num{112000} वर्षों के आवर्तकाल से घूमती है। मान लीजिए कि दोनों गतियाँ
उपसौर को विषुव के सापेक्ष एक ही दिशा में खिसकाती हैं, जिससे उनकी आवृत्तियाँ जुड़ती हैं। (क)~दोनों
आवृत्तियों को मिलाकर वह आवर्तकाल ज्ञात कीजिए जिससे उपसौर
ऋतुओं में से होकर खिसकता है। (ख)~\omterm{def:b2:orbital-forcing:elements}{जलवायवीय पुरस्सरण} के मुख्य आवर्तकालों, \num{19.0},
\num{22.4} और \num{23.7} हज़ार वर्ष, से तुलना कीजिए।
\end{exercise}
% ledger: nasa:milank-detail, la04:precession

\begin{exercise}[$\star$]\label{exo:b2:orbital-forcing:4}
\cref{prop:b2:orbital-forcing:annual} और $e = 0.0167$ से, \ang{22.1} और \ang{24.5} की तिर्यकताओं के लिए
उत्तरी ध्रुव पर वार्षिक-माध्य सूर्यातप परिकलित कीजिए। दोनों के बीच वह कितने प्रतिशत
बढ़ता है?
\end{exercise}
% ledger: nasa:milank-detail, earth:ecc

\begin{exercise}[$\star\star$]\label{exo:b2:orbital-forcing:5}
(क)~दूरी गुणक $(a/r)^2$ को जून अयनांत पर, आज और \num{116000} वर्ष
पहले ($e = 0.043988$, $\varpi = \ang{92.71}$), यथार्थ रूप से परिकलित कीजिए, और सन्निकटन $1 -
2e\sin\varpi$ से तुलना कीजिए। (ख)~दिसंबर अयनांत पर दोनों परिकलित कीजिए। (ग)~\num{116000}
वर्ष पहले सन्निकटन अधिक खराब क्यों है?
\end{exercise}
% ledger: berger91:rows

\begin{exercise}[$\star\star$]\label{exo:b2:orbital-forcing:6}
\cref{ex:b2:orbital-forcing:today} को चरण-दर-चरण पुनरुत्पन्न कीजिए: जून अयनांत पर \ang{65} उ पर
सूर्यास्त का घंटा कोण, दिन की लंबाई, माध्य दूरी पर सूर्यातप और
आज का सूर्यातप।
\end{exercise}
% ledger: berger91:rows

\begin{exercise}[$\star\star$]\label{exo:b2:orbital-forcing:7}
(क)~$\mathrm dt = T(r/a)^2\,\mathrm d\lambda/(2\pi\sqrt{1 - e^2})$ और $(r/a)^2 \approx 1 +
2e\cos(\lambda - \varpi)$ का उपयोग करके दिखाइए कि मार्च से सितंबर विषुव तक का अर्ध-वर्ष
लगभग $T/2 + (2T/\pi)\,e\sin\varpi$ चलता है। (ख)~परिकलित कीजिए कि आज वह दूसरे अर्ध-वर्ष से कितने दिन
अधिक है। (ग)~उन 4.5 दिनों से, जिनसे उत्तरी ग्रीष्म शीत ऋतु से अधिक है, और उन तीन दिनों से, जिनसे
वसंत शरद से अधिक है, तुलना कीजिए।
\end{exercise}
% ledger: nasa:milank-detail, berger91:rows

\begin{exercise}[$\star\star$]\label{exo:b2:orbital-forcing:8}
पिछले दस लाख वर्षों में उत्केंद्रता \num{0.0046} से \num{0.0565} तक रही। (क)~दोनों चरम मानों पर
वैश्विक और वार्षिक माध्य सूर्यातप $S_0/(4\sqrt{1 - e^2})$ परिकलित कीजिए। (ख)~कक्षा ने उसे कितना बदला,
\unit{W.m^{-2}} में और प्रतिशत में? (ग)~\ang{65} उ पर मध्य-जुलाई सूर्यातप के परास,
388 से \qty{497}{W.m^{-2}}, से तुलना कीजिए।
\end{exercise}
% ledger: data:b4-orbital

\begin{exercise}[$\star\star$]\label{exo:b2:orbital-forcing:9}
पश्चता मॉडल और $\tau = \qty{10}{kyr}$ के साथ, 23, 41 और \num{100} हज़ार वर्षों के आवर्तकालों के ज्यावक्रीय प्रणोदन के लिए
हिम-आयतन का सापेक्ष आयाम और पश्चता परिकलित कीजिए। हिमचादर किस पट्टी को
वरीयता देती है, और क्यों?
\end{exercise}

\begin{exercise}[$\star\star\star$]\label{exo:b2:orbital-forcing:10}
\ang{65} उ पर मध्य-जुलाई सूर्यातप में \num{100000} वर्षों पर लगभग कोई शक्ति नहीं है, फिर भी पिछले
\num{700000} वर्षों के हिम-आयतन की अधिकांश शक्ति वहीं है। (क)~दिखाइए कि यदि प्रणोदन $A_1\sin
\omega_1 t + A_2\sin\omega_2 t$ हो, तो पश्चता मॉडल केवल $\omega_1$ और $\omega_2$ पर अनुक्रिया करता है। (ख)~दिखाइए
कि पद $V^2$ वाली अनुक्रिया, उदाहरण के लिए $\mathrm dV/\mathrm dt = (V_{\mathrm{eq}} -
V)/\tau + kV^2$, आवृत्तियाँ $\omega_1 - \omega_2$ और $\omega_1 + \omega_2$ उत्पन्न कर सकती है,
रैखिक हल को $V^2$ में प्रतिस्थापित करके। (ग)~\num{23.7} और
\num{19.0} हज़ार वर्षों के आवर्तकालों वाले दो पुरस्सरण पद किस आवर्तकाल से विस्पंद करते हैं? \num{100000}-वर्षीय
चक्रों के बारे में आप क्या निष्कर्ष निकालते हैं?
\end{exercise}
% ledger: la04:precession, barker21:terminations

\begin{exercise}[$\star\star\star$]\label{exo:b2:orbital-forcing:11}
उत्तरी ध्रुव पर ध्रुवीय दिन के दौरान प्रति वर्ग मीटर प्राप्त ऊर्जा $E = \int S_0(a/r)^2
\sin\delta\,\mathrm dt$ है, उस अर्ध-वर्ष पर जब $\delta > 0$। (क)~\cref{prop:b2:orbital-forcing:annual} की
उपपत्ति से दिखाइए कि $E = S_0T\sin\varepsilon/(\pi\sqrt{1 - e^2})$। (ख)~$E$ को \unit{GJ.m^{-2}} में परिकलित कीजिए, $\varepsilon = \ang{22.1}$ और \ang{24.5} के लिए, $e = 0.0167$ और $T$ =
365.25 दिन के साथ।
(ग)~पुरस्सरण $E$ में प्रकट नहीं होता। फिर भी वह किसी हिमचादर के पिघलने को क्यों
बदल सकता है?
\end{exercise}
% ledger: nasa:milank-detail

\begin{exercise}[$\star\star\star$]\label{exo:b2:orbital-forcing:12}
(क)~कार्बन डाइऑक्साइड के \qty{189.7}{ppm} से \qty{279.4}{ppm} तक बढ़ने का विकिरणी प्रणोदन परिकलित कीजिए,
\qty{3.93}{W.m^{-2}} प्रति दोगुनीकरण और एक लघुगणकीय नियम के साथ। (ख)~अंतिम हिमनदीय अधिकतम और अंतिम \omterm{def:b2:orbital-forcing:glacials}{अंतर्हिमनदीय काल} के लिए
दिए गए परास 188 से \qty{194}{ppm} और 266 से
\qty{282}{ppm} हैं: वे जो सबसे छोटा और सबसे बड़ा प्रणोदन अनुमत करते हैं, उन्हें परिकलित कीजिए। (ग)~\cref{exo:b2:orbital-forcing:8} के
वैश्विक वार्षिक सूर्यातप के परिवर्तन से और \ang{65} उ ग्रीष्म सूर्यातप के उतार-चढ़ाव से
तुलना कीजिए। कौन ट्रिगर है और कौन प्रवर्धक?
\end{exercise}
% ledger: data:b4-proxies-co2, ipcc:fig2-34, ipcc:erf-2xco2, ipcc:log-co2, data:b4-orbital

\begin{problem}[{सप्ताहांत समस्या --- \ang{65} उ पर सूर्य-प्रकाश}]\label{pb:b2:orbital-forcing:1}
अंतिम \omterm{def:b2:orbital-forcing:glacials}{अंतर्हिमनदीय काल} लगभग \num{116000} वर्ष पहले समाप्त हुआ। तब विशाल उत्तरी हिमचादरों के
अक्षांश की ग्रीष्म को आज की तुलना में कितनी धूप मिली? अभ्यासों के आँकड़ों का उपयोग कीजिए।
तालिकाएँ \num{116000} वर्ष पहले के लिए $e = 0.0440$, $\varpi = \ang{92.7}$,
$\varepsilon = \ang{22.52}$ देती हैं; आज के लिए, उसी प्रकार पूर्णांकित, $e = 0.0172$, $\varpi = \ang{101.4}$,
$\varepsilon = \ang{23.45}$।

\medskip\noindent\textbf{भाग I --- तब और अब की कक्षा।}
\begin{enumerate}
  \item तब और अब का \omterm{def:b2:orbital-forcing:elements}{जलवायवीय पुरस्सरण} परिकलित कीजिए।
  \item तब और अब की उपसौर और अपसौर दूरियाँ, $a$ की इकाइयों में, परिकलित कीजिए।
  \item तब, और अब, सूर्य के किस देशांतर पर पृथ्वी उपसौर पर थी? किस अयनांत के निकट?
  \item तब और अब, उपसौर पर सूर्य-प्रकाश अपसौर की तुलना में किस गुणक से अधिक प्रबल था?
  \item क्या तिर्यकता आज से अधिक थी या कम? उच्च अक्षांशों की ग्रीष्म ऋतुओं के लिए
        इसका क्या अर्थ है?
  \item तब और अब का वैश्विक और वार्षिक माध्य सूर्यातप परिकलित कीजिए। क्या अंतर महत्त्वपूर्ण है?
\end{enumerate}

\medskip\noindent\textbf{भाग II --- जून अयनांत पर दूरी।}
\begin{enumerate}[resume]
  \item \cref{prop:b2:orbital-forcing:distance} से दिखाइए कि जून अयनांत पर $r/a = (1 -
        e^2)/(1 - e\sin\varpi)$।
  \item तब और अब, जून अयनांत पर $r/a$ परिकलित कीजिए।
  \item तब और अब, $(a/r)^2$ परिकलित कीजिए।
  \item सन्निकटन $1 - 2e\sin\varpi$ से तुलना कीजिए।
  \item तब और अब, अयनांत का सूर्य-प्रकाश माध्य दूरी की तुलना में कितने प्रतिशत दुर्बल था?
\end{enumerate}

\medskip\noindent\textbf{भाग III --- \ang{65} उ पर धूप।}
\begin{enumerate}[resume]
  \item तब और अब, जून अयनांत पर सूर्य की क्रांति दीजिए।
  \item तब और अब, उस दिन \ang{65} उ पर सूर्यास्त का घंटा कोण परिकलित कीजिए।
  \item दिन की लंबाई निकालिए।
  \item माध्य दूरी पर दैनिक-माध्य सूर्यातप परिकलित कीजिए।
  \item तब और अब, दैनिक-माध्य सूर्यातप परिकलित कीजिए।
  \item तालिकाएँ मध्य-जुलाई के लिए 397.60 और \qty{426.76}{W.m^{-2}} देती हैं। अयनांत के
        मान अधिक क्यों हैं?
  \item तब और अब के अंतर में किसका योगदान अधिक था, तिर्यकता का या दूरी का?
\end{enumerate}

\medskip\noindent\textbf{भाग IV --- हिमचादर का उत्तर।}
\begin{enumerate}[resume]
  \item $\tau = \qty{10}{kyr}$ वाले पश्चता मॉडल में, हिमचादर
        $V_{\mathrm{eq}}$ के एक आकस्मिक परिवर्तन का कितना अंश 5, 10 और \qty{20}{kyr} के बाद प्राप्त कर चुकी होती है?
  \item \ang{65} उ पर जून सूर्यातप लगभग \num{128000} वर्ष पहले के एक अधिकतम से
        \num{116000} वर्ष पहले के न्यूनतम तक गिरा, यानी आधा पुरस्सरण-चक्र। आवर्तकाल
        \qty{23}{kyr} के ज्यावक्र के लिए हिमचादर प्रणोदन के कितने अंश का अनुसरण करती है, और कितनी देर से?
  \item कोई हिमचादर एक पुरस्सरण-चक्र के भीतर वह आकार क्यों नहीं पा सकती जिसे एक निम्न ग्रीष्म
        सूर्यातप बनाए रखेगा?
  \item बर्फ़ को अंतिम हिमनदीय अधिकतम तक बनने में लगभग \num{90000} वर्ष लगे और पिघलने में लगभग
        \num{10000} वर्ष। पश्चता मॉडल यह विषमता क्यों नहीं उत्पन्न कर सकता?
  \item अंतिम \omterm{def:b2:orbital-forcing:glacials}{अंतर्हिमनदीय काल} में कार्बन डाइऑक्साइड 266 और \qty{282}{ppm} के बीच थी और अंतिम हिमनदीय अधिकतम पर
        188 और \qty{194}{ppm} के बीच। वह धूप के एक क्षेत्रीय परिवर्तन को जलवायु के एक वैश्विक परिवर्तन
        बनने में कैसे सहायता करती है?
  \item सारांश में: जून अयनांत पर \ang{65} उ पर दैनिक-माध्य सूर्यातप कितना था, जब
        अंतिम \omterm{def:b2:orbital-forcing:glacials}{अंतर्हिमनदीय काल} समाप्त हुआ, \num{116000} वर्ष पहले?
\end{enumerate}
\end{problem}
% ledger: berger91:rows, berger91:readme, data:b4-orbital, ipcc:ref-periods, ipcc:fig2-34, barker21:terminations
